%====================================================================
%  Stage 1: Fixed Directed-Cost Spatial Equilibrium
%  Beijing, raw_connected v1 edgeproj_v2, LCC + detour gate G=5
%====================================================================
\documentclass[12pt]{article}

\usepackage[margin=1in]{geometry}
\usepackage{amsmath,amssymb,amsthm,mathtools}
\usepackage{graphicx,booktabs,tabularx,array,setspace,float,enumitem}
\usepackage{natbib}
\usepackage[colorlinks=true,linkcolor=blue,citecolor=blue,urlcolor=blue]{hyperref}
\setstretch{1.15}
\setlength{\parindent}{1.5em}

\newtheorem{definition}{Definition}
\newtheorem{assumption}{Assumption}
\newtheorem{remark}{Remark}

\newcommand{\N}{\mathcal N}
\newcommand{\E}{\mathcal E}
\newcommand{\Lbar}{\bar L}
\newcommand{\Wbar}{\bar W}
\newcommand{\ubar}{\bar u}
\newcommand{\Abar}{\bar A}
\newcommand{\obs}{\mathrm{obs}}
\newcommand{\sym}{\mathrm{sym}}
\newcommand{\dir}{\mathrm{dir}}

\graphicspath{{../../data_work/outputs/directional_cost_edgeproj_v2_lcc_gate5_v1/figures/}}

\title{Stage~1: Fixed Directed-Cost Spatial Equilibrium\\[0.4em]
\large Beijing Morning Peak, \texttt{model\_input\_edgeproj\_v2\_lcc\_gate5\_v1}}
\author{Yige Duan \and Xiaoruo Feng \and Yizhen Gu}
\date{\today}

\begin{document}
\maketitle

\begin{abstract}
\noindent
This note documents the first structural exercise in the paper's two-stage model line.
Observed directed peak-period travel times are held fixed within each exercise; endogenous
congestion feedback is excluded throughout.
For each cost surface (directed observed vs.\ symmetric approximation), the code
recovers baseline residential amenities and workplace productivities that rationalise
the observed road-commuter marginals, then solves the static spatial equilibrium.
We further run non-directional transport counterfactuals, a tidal-lane \emph{prototype}
(no congestion spillovers), and asymmetry sensitivity grids.
The central question is \emph{which observable planning objects}---commuting flows,
residential and workplace sorting, and grid-level job access---depend on retaining
directional peak costs rather than a symmetric link index.
Within-run welfare $\Wbar$ moves with inverted fundamentals and is reported for
completeness, but it is not the headline object: under re-inversion both worlds fit
the same marginals, so welfare differences largely reflect compensation in
$(\bar u,\bar A)$ rather than a policy-relevant quantity planners observe directly.
\end{abstract}

\section{Research question and scope}\label{sec:question}

Directional congestion is prevalent in Beijing's morning peak.
A symmetric travel-cost index averages the two directions of a corridor and therefore
mixes a slow peak direction with a fast off-peak direction on the same physical link.
Stage~1 asks:
\begin{quote}
\noindent
\emph{If we build a directed calibration world and a symmetric calibration world---each
with its own inverted fundamentals and each fitting the same observed residence and
workplace marginals---which \textbf{observable outcomes} differ, and which classes of
transport or place-based policy are mis-measured when directional peak congestion is
collapsed to a symmetric link index?}
\end{quote}

\paragraph{Headline objects vs.\ welfare.}
Planners observe population counts, commute patterns, and accessibility by zone---not
within-run $\Wbar$.
Because inversion pins $(l_i^R,l_j^F)$ to the data in each world, baseline sorting
across worlds is negligible and baseline $\Wbar$ gaps ($\sim$1\%) mainly trace
re-inverted fundamentals.
We therefore emphasise (i)~predicted OD structure, (ii)~grid-level market access
$\Pi_i^{-\theta}$, and (iii)~how \emph{post-calibration} transport shocks move
residence shares, hub-bound commuting, and access inequality---objects that map to
school and clinic catchments, transit ridership to employment centres, and corridor
investment priorities.

Three objects remain \emph{out of scope} for Stage~1:
(i)~endogenous link traffic and congestion feedback (including spillovers from tidal lanes
to untreated edges);
(ii)~structural estimation of $(\theta,\alpha,\beta)$ or congestion elasticity $\lambda$;
(iii)~full tidal-lane welfare evaluation (only a fixed-$\xi$ BPR prototype is reported in
\S\ref{sec:tidal_lane}).

\section{Environment and data}\label{sec:data}

\subsection{Spatial units and network scope}

The city is partitioned into $N=1104$ square grid cells in the largest weakly connected
component of the directed adjacency graph implied by pair-specific routed travel costs.
The retained network contains $|\E|=2642$ directed edges and $310{,}525$ origin--destination
pairs with positive road-mode commuters ($\Lbar=88{,}110.2$ road commuters; $92.6\%$ of the
city-wide road-commuter total in the baseline OD table).

Travel costs come from \texttt{raw\_connected\_v1} edge-projection routing
(\texttt{edgeproj\_v2}) with detour gate $G=5$ and LCC restriction.
Input QA is in \path{model_input_qa.json}.

\subsection{Observed objects}

Each directed edge $(k,l)\in\E$ carries observed peak-period travel time
$\tau^{\obs}_{kl}>0$ (minutes).
There are $1193$ unordered reciprocal pairs; $72.4\%$ have peak-period travel-time ratio
$r_{kl}\equiv\max(\tau^{\obs}_{kl},\tau^{\obs}_{lk})/\min(\cdot)\ge 1.10$.

Road-mode commuting flows $\{L^{\obs}_{ij}\}$ yield marginals
$L^{\obs,R}_i=\sum_j L^{\obs}_{ij}$ and $L^{\obs,F}_j=\sum_i L^{\obs}_{ij}$.
Period: AM; population column: \texttt{commuters\_road}.

\section{Model}\label{sec:model}

\subsection{Fixed travel cost}

Stage~1 holds the edge cost vector fixed within each world:
\begin{equation}
  \tau_{kl} \;=\; \tau^{\obs}_{kl}
  \quad\text{(directed world)},\qquad
  \tau_{kl} \;=\; \tau^{\sym}_{kl}
  \quad\text{(symmetric world)}.
  \label{eq:fixed_tau}
\end{equation}
No link updates from lane counts, traffic, or congestion feedback.

\subsection{Route aggregator}

\begin{equation}
  A_{kl} \;=\; \frac{1}{\kappa\,\tau_{kl}},
  \qquad
  T_{ij} \;\equiv\; \bigl[(I-A)^{-1}\bigr]_{ij},
  \qquad
  \tau_{ij}^{-\theta} \;\equiv\; T_{ij},
  \label{eq:route_aggregator}
\end{equation}
with Allen--Arkolakis iceberg convention and Fréchet shape $\theta$
\citep{AllenArkolakis2022}.
On the gated LCC network, raw row sums of $1/\tau_{kl}$ are already below $0.9$, so
$\kappa=1$.

\subsection{Fundamentals and static equilibrium}

Baseline objects $\ubar_i>0$ and $\Abar_j>0$ enter with agglomeration:
$u_i=\ubar_i(L_i^R)^{\beta}$, $A_j=\Abar_j(L_j^F)^{\alpha}$.
Shares $(l_i^R,l_j^F)$ solve
\begin{align}
  l_i^R
  &= \left[
      \chi\,\ubar_i^{\theta}\,
      \sum_j T_{ij}\,
      \Abar_j^{\theta}\,
      (l_j^F)^{\alpha\theta}
    \right]^{\!1/(1-\beta\theta)},
  \label{eq:eq_R} \\[0.4em]
  l_j^F
  &= \left[
      \chi\,\Abar_j^{\theta}\,
      \sum_i T_{ij}\,
      \ubar_i^{\theta}\,
      (l_i^R)^{\beta\theta}
    \right]^{\!1/(1-\alpha\theta)},
  \label{eq:eq_F}
\end{align}
with $\sum_i l_i^R=\sum_j l_j^F=1$.
Within-run welfare is
\begin{equation}
  \Wbar \;\equiv\;
  \left[
    \sum_i \ubar_i^{\theta}(L_i^R)^{\beta\theta}\,
    \Pi_i^{-\theta}
  \right]^{1/\theta}.
  \label{eq:welfare}
\end{equation}

\subsection{Fundamentals inversion}\label{sec:inversion}

Given a fixed cost surface $\{\tau_{kl}\}$ and observed shares
$(l_i^{R,\obs},l_j^{F,\obs})$, inversion finds
$(\ubar_i^{\theta},\Abar_j^{\theta},\chi)$ such that equations
\eqref{eq:eq_R}--\eqref{eq:eq_F} hold at the observed marginals.
Equivalently, with $v_i\equiv\ubar_i^{\theta}$, $w_j\equiv\Abar_j^{\theta}$,
$p_R=1-\beta\theta$, $p_F=1-\alpha\theta$:
\begin{align}
  v_i \bigl(T (w \odot (l^{F,\obs})^{\alpha\theta})\bigr)_i
  &= \chi\,(l_i^{R,\obs})^{p_R},
  \label{eq:invert_R}\\
  w_j \bigl(T^{\top} (v \odot (l^{R,\obs})^{\beta\theta})\bigr)_j
  &= \chi\,(l_j^{F,\obs})^{p_F}.
  \label{eq:invert_F}
\end{align}
The algorithm alternates updates on \eqref{eq:invert_R}--\eqref{eq:invert_F}, and after
each half-step rescales $v$ and $w$ to geometric mean one on valid nodes
($L^{\obs,R}_i,L^{\obs,F}_i>0$), implementing
$\prod_i\ubar_i=\prod_j\Abar_j=1$.
Convergence requires the implied node-level $\chi$ from the residence and workplace
equations to agree (std$(\log\chi)<1.5\times 10^{-4}$).
This is a structural fixed point, not a gravity regression with origin/destination fixed
effects.

\subsection{Symmetric-cost world}

For reciprocal pairs,
\begin{equation}
  \tau^{\sym}_{kl}=\tau^{\sym}_{lk}
  \;=\;
  \sqrt{\tau^{\obs}_{kl}\,\tau^{\obs}_{lk}}
  \qquad\text{(geometric mean; main specification)}.
  \label{eq:sym_geometric}
\end{equation}
One-way edges keep $\tau^{\obs}_{kl}$.
The symmetric world \emph{re-inverts} $(\ubar,\Abar,\chi)$ under $\tau^{\sym}$ and
re-solves \eqref{eq:eq_R}--\eqref{eq:eq_F}.
Both worlds target the same observed marginals; differences arise in OD flows, welfare
levels, and the inverted baseline objects.

\begin{remark}[Link-level symmetrisation, not OD symmetrisation]
Symmetrisation applies to \emph{directed edges} on the routing graph: for each unordered
node pair $\{k,l\}$ with reciprocal links, both directions receive the same
$\tau^{\sym}_{kl}=\tau^{\sym}_{lk}$.
It does \emph{not} symmetrise the observed OD matrix $L^{\obs}_{ij}$ or replace
home-to-work costs by work-to-home costs.
AM peak OD is home $\to$ work; edge costs are directed AM travel times $k\to l$.
QA (\S\ref{sec:qa_direction}) confirms that shortest-path costs on $\tau^{\obs}$ align
with this convention for high-flow off-diagonal pairs.
\end{remark}

\subsection{Implied OD flows (diagnostic, not inverted)}\label{sec:od_gravity}

Given equilibrium shares and inverted fundamentals, bilateral flows follow a gravity
allocation on the Leontief aggregator (not targeted in inversion):
\begin{equation}
  L_{ij}
  \;=\;
  L_i^R\,
  \frac{
    T_{ij}\,\Abar_j^{\theta}\,(l_j^F)^{\alpha\theta}
  }{
    \sum_k T_{ik}\,\Abar_k^{\theta}\,(l_k^F)^{\alpha\theta}
  },
  \qquad L_i^R = \Lbar\, l_i^R.
  \label{eq:od_gravity}
\end{equation}
Dissimilarity between two flow matrices is reported as
$0.5\sum_{i,j}|L_{ij}^{(1)}-L_{ij}^{(2)}|/\sum_{i,j}L_{ij}^{(1)}$ in percentage units.
Stage~1 compares implied OD across calibration worlds; fit to $L^{\obs}_{ij}$ is
documented but not optimised (\S\ref{sec:od_fit}).

\subsection{Tidal-lane travel-time update (prototype)}\label{sec:tidal_bpr}

For a reciprocal pair with slow direction $s$ and fast direction $f$, lane reallocation
$\Delta n>0$ moves capacity from $f$ to $s$ without updating congestion on other links.
With BPR exponent $\delta_1=0.488$ and fixed observed congestion index
$\xi_s = n_s(\max\{\tau^{\obs}_s,\tau^{ff}_s\}/\tau^{ff}_s)^{1/\delta_1}$:
\begin{equation}
  \tau_s^{\mathrm{new}}
  \;=\;
  \max\left\{
    \tau^{ff}_s,\;
    \tau^{ff}_s\left(\frac{\xi_s}{n_s+\Delta n}\right)^{\delta_1}
  \right\},
  \qquad
  \tau_f^{\mathrm{new}}
  \;=\;
  \max\left\{
    \tau^{ff}_f,\;
    \tau^{ff}_f\left(\frac{\xi_f}{n_f-\Delta n}\right)^{\delta_1}
  \right\}.
  \label{eq:tidal_bpr}
\end{equation}
Only edges in the selected pair set change; all other $\tau_{kl}$ remain at
$\tau^{\obs}_{kl}$.
Short-run evaluation fixes $(l_i^R,l_j^F)$ and recomputes $\Pi_i$, $\Wbar$, and
\eqref{eq:od_gravity}; long-run re-solves \eqref{eq:eq_R}--\eqref{eq:eq_F}.
Details and grids: \S\ref{sec:tidal_lane}.

\section{Computation and reproducibility}\label{sec:computation}

\subsection{Pipeline overview}

Main runner (\texttt{run\_directional\_cost.py}):
\begin{enumerate}[leftmargin=2em,itemsep=0.2em]
  \item Load QSM primitives; floor zero-mass shares at $10^{-6}$.
  \item Build $T$ from \eqref{eq:route_aggregator}; set $\kappa=1$.
  \item \textbf{Directed world:} invert under $\tau^{\obs}$; solve equilibrium.
  \item \textbf{Symmetric world:} build $\tau^{\sym}$; invert again; solve equilibrium.
  \item Report comparisons; asymmetry intensity/threshold grids (re-invert at each surface).
  \item Optional OD-fit diagnostics (\texttt{od\_fit\_estimation.py}).
\end{enumerate}

Additional runners (same input bundle, directed re-invert baseline unless noted):
\begin{center}
\small
\begin{tabular}{lll}
\toprule
Script & Purpose & Output subfolder \\
\midrule
\texttt{run\_policy\_counterfactuals.py} & 5 non-directional CFs $\times$ 2 worlds
  & \texttt{counterfactuals/} \\
\texttt{run\_tidal\_lane\_counterfactuals.py} & Tidal-lane grid, short vs long
  & \texttt{counterfactuals/tidal\_lane/} \\
\texttt{plot\_results.py} & Asymmetry histograms, simulation curves
  & \texttt{figures/} \\
\texttt{plot\_stage1\_maps.py} & $\Pi_i^{-\theta}$, $P_j^{-\theta}$, wage maps
  & \texttt{figures/} \\
\texttt{plot\_tidal\_lane\_results.py} & Short vs long tidal-lane scatter
  & \texttt{figures/} \\
\texttt{qa\_od\_direction\_vs\_cost.py} & AM OD direction vs $\tau$ QA
  & \texttt{diagnostics/od\_direction\_cost\_qa.json} \\
\texttt{plot\_model\_mechanics.py} & Obs vs model workplace-choice bars
  & external or \texttt{figures/} \\
\bottomrule
\end{tabular}
\end{center}

All commands assume repository root \path{TidalLanes/}:
\begin{verbatim}
python3 data_work/src/model/directional_cost/run_directional_cost.py \
  --version-id model_input_edgeproj_v2_lcc_gate5_v1 \
  --output-name directional_cost_edgeproj_v2_lcc_gate5_v1
python3 data_work/src/model/directional_cost/run_policy_counterfactuals.py
python3 data_work/src/model/directional_cost/run_tidal_lane_counterfactuals.py
python3 data_work/src/model/directional_cost/plot_stage1_maps.py \
  --output-root data_work/outputs/directional_cost_edgeproj_v2_lcc_gate5_v1
\end{verbatim}

\subsection{Convergence and numerical settings}

Directed inversion: $503$ iterations, $\max|\Delta\log\chi|=1.49\times 10^{-4}$.
Symmetric inversion: $511$ iterations, same tolerance
(std$(\log\chi)<1.5\times 10^{-4}$).
Static equilibria converge in $10$ iterations each
($\max|\Delta\log l|<10^{-9}$).
Inversion tolerance, static tolerance, and intensity/threshold grids are stored in
\path{run_metadata.json}.

\subsection{Output manifest}\label{sec:output_manifest}

Output root:
\texttt{data\_work/outputs/directional\_cost\_edgeproj\_v2\_lcc\_gate5\_v1/}.
Key artifacts:
\begin{itemize}[leftmargin=2em,itemsep=0.15em]
  \item \path{summary_metrics.json} --- directed vs symmetric headline comparison.
  \item \path{fundamentals/directed|symmetric/} --- inverted $\ubar_i^{\theta}$,
        $\Abar_j^{\theta}$, $\chi$.
  \item \path{equilibrium/directed_observed|symmetric_cost/} --- $l_i^R,l_j^F,\Pi_i,P_j$.
  \item \path{simulations/asymmetry_{intensity,threshold}.csv} --- $\rho$ and $c$ grids.
  \item \path{diagnostics/} --- edge/node comparisons, OD fit JSON, MA map CSV.
  \item \path{counterfactuals/policy_cf_{summary,world_wedge}.csv}.
  \item \path{counterfactuals/tidal_lane/} --- 26-spec grid, pair ranks, edge-level $\tau$ shifts.
  \item \path{figures/} --- all plots referenced in this note.
\end{itemize}
Extended workflow log:
\path{Documents/L2_notes/stage1_fixed_directed_cost_complete_handoff_20260628.md}.

\subsection{Design choices locked for Stage~1}\label{sec:design_locks}

\begin{center}
\small
\begin{tabular}{p{0.28\textwidth}p{0.62\textwidth}}
\toprule
Choice & Stage~1 lock-in \\
\midrule
Input bundle & \texttt{model\_input\_edgeproj\_v2\_lcc\_gate5\_v1} (LCC, detour $G=5$) \\
Cost primitive & $\tau_{kl}=\tau^{\obs}_{kl}$ (AM peak, minutes) \\
Route aggregator & $T=(I-A)^{-1}$, $A_{kl}=1/(\kappa\tau_{kl})$, $\kappa=1$ \\
Symmetrisation & Geometric mean on reciprocal \emph{links} (main); arithmetic robustness \\
Inversion & Structural fixed point on marginals; not gravity FE \\
Comparison mode & Re-invert fundamentals in \emph{both} worlds (main) \\
Parameters & $\theta=6.83$, $\alpha=-0.12$, $\beta=-0.10$ (AA Seattle placeholders) \\
Congestion & None; $\lambda$ omitted \\
Policy CFs & Fixed inverted fundamentals; levels, not exact-hat \\
Tidal lanes & BPR \eqref{eq:tidal_bpr} at fixed $\xi$; no spillovers \\
\bottomrule
\end{tabular}
\end{center}

\section{Parameters and levers that amplify world differences}\label{sec:parameters}

Stage~1 does not estimate $(\theta,\alpha,\beta)$; placeholders from AA Seattle:
$\theta=6.83$, $\alpha=-0.12$, $\beta=-0.10$.
Congestion elasticity $\lambda$ is omitted.

\subsection{What makes the two worlds look similar?}

Under the main \textbf{re-inversion} design, both calibration worlds target the same
observed residence and workplace marginals.
Inversion therefore absorbs much of the symmetric-cost bias into
$(\bar u_i,\bar A_j)$: residence and workplace \emph{shares} differ by less than
$0.001\%$ at baseline, and predicted OD dissimilarity is stable at $\approx 2.5\%$
across a wide range of $\theta$ (Table~\ref{tab:theta}).
This is intentional---it isolates \emph{measurement of commuting technology} from
\emph{re-fitting population counts}---but it also compresses baseline sorting gaps.

\subsection{Knobs that widen directed vs.\ symmetric gaps}

Table~\ref{tab:theta} and the discussion below summarise practical levers.
None changes the Stage~1 scope (fixed $\tau$, no congestion); they alter how large
the cross-world wedge appears on different objects.

\begin{table}[htbp]
  \centering
  \caption{Sensitivity of cross-world gaps to $\theta$ (geometric symmetrisation, re-invert)}
  \label{tab:theta}
  \begin{tabular}{lrrrr}
  \toprule
  $\theta$ & $\Wbar$ change (\%) & OD dissim.\ (\%) & Res.\ dissim.\ (\%) & p95 $|\Delta\log\Pi_i^{-\theta}|$ \\
  \midrule
  3   & $-2.58$ & 2.46 & $<0.001$ & 0.52 \\
  4   & $-1.94$ & 2.46 & $<0.001$ & 0.52 \\
  6.83 (baseline) & $-1.14$ & 2.46 & $<0.001$ & 0.52 \\
  9   & $-0.87$ & 2.46 & $<0.001$ & 0.52 \\
  12  & $+0.27$ & 2.38 & 0.69 & 0.48 \\
  \bottomrule
  \end{tabular}
  \vspace{0.4em}

  \begin{minipage}{0.92\textwidth}
  \footnotesize
  \emph{Notes:} Symmetric minus directed; each $\theta$ re-inverts fundamentals in both
  worlds. OD and $\Pi^{- \theta}$ gaps are robust; $\Wbar$ and residence sorting respond
  more to $\theta$; very high $\theta$ begins to move baseline shares because the fixed
  point is stiffer.
  \end{minipage}
\end{table}

\begin{enumerate}[leftmargin=2em,itemsep=0.35em]
  \item \textbf{Hold fundamentals fixed (legacy contrast).}
        Invert only on the directed surface, then compare directed vs.\ symmetric
        $\tau$ \emph{without} re-inverting in the symmetric world.
        Sorting and OD gaps widen to $\approx 2.9\%$ / $2.5\%$ / $5.6\%$ (residence /
        workplace / OD dissimilarity) because symmetric costs can no longer compensate
        through $(\bar u,\bar A)$.
        Useful as an upper bound on spatial re-sorting from mis-measured $T$, but not the
        main specification because it mixes cost bias with a different fundamental surface.

  \item \textbf{Raise $\theta$ (moderate effect on sorting under re-invert).}
        Lower $\theta$ magnifies the $\Wbar$ gap ($-2.6\%$ at $\theta=3$) but barely
        moves baseline OD or $\Pi^{- \theta}$ maps.
        Policy CFs that load on $T$ (hub access) show slightly larger \emph{post-shock}
        cross-world wedges in upper-tail job access (p95 $\Pi^{- \theta}$ gap $+3.7\%$
        after hub-access CF vs.\ $-4.0\%$ at baseline).

  \item \textbf{Asymmetry intensity / threshold grids} (\S\ref{sec:simulations}).
        Partially restore directed edges ($\rho\in(0,1)$ or threshold $c>1$) smoothly
        interpolates between worlds and shows that even ``moderate'' symmetrisation
        ($c=2$, 8.9\% of reciprocal pairs remain directed) leaves a $0.86\%$ $\Wbar$
        gap relative to full symmetrisation---again mostly fundamental compensation,
        but useful for robustness.

  \item \textbf{Concentrate shocks on high-asymmetry corridors (Stage~2 preview).}
        Non-directional CFs in Stage~1 apply the same scalar speed factor to both
        directions of a link.
        Directional capacity reallocation (tidal lanes) will load on edges with large
        $r_{kl}=\max(\tau_{kl},\tau_{lk})/\min(\cdot)$; Stage~1's hub/fringe transport
        CFs are conservative proxies for \emph{where} mis-measurement bites, not its
        upper bound.

  \item \textbf{Relax marginal pinning (hypothetical).}
        Dropping the re-inversion step---or targeting OD moments in addition to
        marginals---would allow bilateral misfit and asymmetric costs to feed back into
        $(l_i^R,l_j^F)$; that is the Stage~2 / full-estimation agenda, not Stage~1.
\end{enumerate}

\section{Results}\label{sec:results}

\subsection{Directional asymmetry}

Figure~\ref{fig:asym_hist}: median $r=1.23$; p90 $=1.94$; $8.89\%$ of reciprocal pairs
have $r\ge 2.00$.

\begin{figure}[htbp]
  \centering
  \includegraphics[width=0.78\textwidth]{asymmetry_ratio_histogram.png}
  \caption{Reciprocal-pair travel-time ratios.}
  \label{fig:asym_hist}
\end{figure}

\subsection{Fit to observed marginals}

Both calibration worlds reproduce floored observed residence and workplace shares:
L1 distance $\approx 0.00053\%$ (directed) and $\approx 0.00053\%$ (symmetric).
Inversion is tied to \emph{marginals}, not the full OD matrix.

\subsection{Implied OD vs.\ observed OD (diagnostic scope)}\label{sec:od_fit}

Stage~1 does \emph{not} target bilateral fit.
Given inverted fundamentals and the gravity allocation
$L_{ij}\propto L_i^R T_{ij}\bar A_j^{\theta}(l_j^F)^{\alpha\theta}$,
the directed calibration world yields a global L1 distance to observed OD flows of
$\approx 72\%$ on positive cells ($\approx 51\%$ off-diagonal).
Residence-conditional destination shares have median L1 $\approx 66\%$.
The largest gaps concentrate on \emph{intra-cell} flows ($i=j$): the model assigns
$\approx 53\%$ of predicted commuters within the same grid cell versus $\approx 8\%$
in the data, because the Leontief diagonal $T_{ii}$ is large relative to off-diagonal
$T_{ij}$.
This poor bilateral fit does \emph{not} invalidate the Stage~1 question---which compares
two calibration worlds that both fit the same marginals---but it rules out interpreting
Stage~1 as a structural replication of the Beijing OD matrix.
Symmetric vs.\ directed cost changes \emph{model-implied} OD by only $\approx 2.5\%$
while leaving fit to observed OD essentially unchanged ($\approx 72\%$ L1 either way).
Full diagnostics and improvement priorities are logged in
\path{Documents/L2_notes/stage1_od_fit_diagnosis_and_handoff_20260628.md}.

\subsection{AM direction and travel-cost QA}\label{sec:qa_direction}

Script \texttt{qa\_od\_direction\_vs\_cost.py} checks that AM OD flows $i\to j$
(home to work) are evaluated against directed peak costs consistent with edge
$\tau^{\obs}_{k\to l}$.
On off-diagonal positive-flow pairs with finite shortest paths in both directions,
$57.4\%$ have home-to-work SP cost greater than work-to-home SP cost (median log gap
$+0.075$ in minutes-log units), confirming measurable morning peak directionality at the
OD level.
Among reachable off-diagonal pairs, weighted log correlation between SP$(i,j)$ and
soft model cost $T_{ij}^{-1/\theta}$ is $0.94$ (flow-weighted subsample: $0.93$).
This QA supports using directed AM $\tau$ for workplace choice while documenting that
Leontief $T_{ij}$ is a smoother network aggregator than raw shortest-path costs on
sparse high-flow cells.
Full JSON: \path{diagnostics/od_direction_cost_qa.json}.

\subsection{Main comparison: symmetric vs.\ directed (geometric symmetrisation)}

Table~\ref{tab:main}.
With re-inversion, residence and workplace \emph{shares} barely move (dissimilarity
$<0.001\%$): both worlds rationalise the same observed marginals.
Predicted \emph{OD flows} still differ by $2.46\%$ dissimilarity because bilateral
route costs and inverted $\Abar_j$ differ.
Within-run $\Wbar$ falls by $1.14\%$ in the symmetric world.
Fundamentals shift materially: p95 $|\log(\ubar_i^{\sym}/\ubar_i^{\dir})|\approx 0.55$
(and similarly for $\Abar$).

\begin{table}[htbp]
  \centering
  \caption{Symmetric geometric world relative to directed world (both re-inverted)}
  \label{tab:main}
  \begin{tabular}{lc}
  \toprule
  Object & Value \\
  \midrule
  $\Wbar$ change (\%) & $-1.141$ \\
  Residence-share dissimilarity (\%) & $<0.001$ \\
  Workplace-share dissimilarity (\%) & $<0.001$ \\
  Predicted OD dissimilarity (\%) & $2.463$ \\
  Mean edge $\tau^{\sym}/\tau^{\obs}-1$ (\%) & $+1.869$ \\
  OD-weighted bilateral route cost change (\%) & $-1.028$ \\
  p95 $|\log(\ubar_i^{\sym}/\ubar_i^{\dir})|$ (valid nodes) & $0.550$ \\
  p95 $|\log(\Abar_j^{\sym}/\Abar_j^{\dir})|$ (valid nodes) & $0.554$ \\
  \bottomrule
  \end{tabular}
  \vspace{0.4em}

  \begin{minipage}{0.92\textwidth}
  \footnotesize
  \emph{Notes:} Dissimilarity is $0.5\sum|\Delta\text{share}|$.
  Comparison is between two self-consistent calibration worlds, not a fixed-fundamentals
  counterfactual. Legacy fixed-fundamentals mode (invert on directed only): residence /
  workplace / OD dissimilarity $\approx 2.9\% / 2.5\% / 5.6\%$.
  \end{minipage}
\end{table}

Arithmetic symmetric operator (re-inverted): $\Wbar$ change $-0.85\%$; OD dissimilarity
$2.29\%$ (\texttt{directional\_cost\_edgeproj\_v2\_lcc\_gate5\_v1\_sym\_arithmetic/}).

\subsection{Market access, firm access, and wage indices}\label{sec:ma_maps}

Even when residence and workplace \emph{shares} barely move across calibration worlds,
market-access indices and inverted productivities differ spatially.
Let $\Pi_i^{-\theta}$ denote residence-side job access and $P_j^{-\theta}$ firm-side
labor-market access (equilibrium objects from \eqref{eq:welfare}).
The effective productivity index is $A_j=\bar A_j(L_j^F)^{\alpha}$ evaluated at each
world's inverted fundamentals and equilibrium shares.

Figure~\ref{fig:ma_maps}: grid-level maps (square-grid index coordinates).
Symmetric minus directed differences reach p95 $|\Delta\log \Pi_i^{-\theta}|\approx 0.52$
and p95 $|\Delta\log P_j^{-\theta}|\approx 0.53$; the wage index $A_j$ differs mainly
through re-inverted $\bar A_j$ (p95 $|\Delta\log A_j|\approx 0.08$) rather than through
$L_j^F$, which is pinned by observed marginals.

\begin{figure}[htbp]
  \centering
  \includegraphics[width=\linewidth]{map_residence_access.png}\\[0.4em]
  \includegraphics[width=\linewidth]{map_firm_access.png}\\[0.4em]
  \includegraphics[width=\linewidth]{map_wage_productivity.png}
  \caption{Grid maps: residence access $\Pi_i^{-\theta}$, firm access $P_j^{-\theta}$,
  and $A_j=\bar A_j(L_j^F)^{\alpha}$ in the directed world, the symmetric world, and
  symmetric minus directed.}
  \label{fig:ma_maps}
\end{figure}

\subsection{Transmission from calibration-world bias to policy analysis}\label{sec:transmission}

Symmetrising link costs affects policy analysis through four linked channels.
The object of interest is how an \emph{exogenous} post-calibration shock propagates
differently depending on which world supplied $(\tau,T,\bar u,\bar A,\Pi,P)$, and
which \textbf{planner-observable outcomes} move as a result.

\paragraph{Channel 1: Mis-measured commuting technology.}
Replacing $\tau^{\obs}_{kl}$ by a pair average changes the Leontief aggregator
$T_{ij}=[(I-A)^{-1}]_{ij}$ and hence residence-side job access $\Pi_i^{-\theta}$
and firm-side labor access $P_j^{-\theta}$.
This channel is active whenever a policy shock modifies travel costs or re-routes
commuters through $T_{ij}$---i.e.\ road speed-ups, hub accessibility investments,
outer-ring connector upgrades, and (in Stage~2) directional lane reallocation.

\paragraph{Channel 2: Fundamental absorption at inversion.}
Because inversion targets observed marginals only, symmetric costs shift inverted
$\bar u_i$ and $\bar A_j$ (p95 $|\log(\bar u_i^{\sym}/\bar u_i^{\dir})|\approx 0.55$).
The symmetric world starts from a different compensated amenity/productivity surface
even though $l_i^R$ and $l_j^F$ match the data.
This channel matters for interpreting \emph{levels} of $\Wbar$ and inverted
fundamentals, but not for policies that only scale local $\bar A_j$ or $\bar u_i$ on
fixed geographic sets once those fundamentals are held fixed in the CF.

\paragraph{Channel 3: General-equilibrium propagation.}
An exogenous shock is applied on top of \emph{that world's} inverted fundamentals.
The same shock rule hits different $(T,\Pi,P,\bar u,\bar A)$ baselines and generates
different $\Delta l_i^R$, $\Delta l_j^F$, $\Delta L_{ij}$, and $\Delta\Pi_i^{-\theta}$.

\paragraph{Channel 4: Sorting vs.\ commuting structure.}
Baseline sorting differences across worlds are tiny under re-inversion.
Transport policy shocks still re-sort population and reallocate OD flows
(dissimilarity $\sim$2--7\% within each world), and cross-world \emph{post-shock}
gaps in upper-tail job access widen even when population-share gaps stay below
$0.02$ percentage points.

\medskip
\noindent\textbf{Policy taxonomy (Stage~1).}
\begin{center}
\small
\begin{tabular}{p{0.26\textwidth}p{0.30\textwidth}p{0.36\textwidth}}
\toprule
Shock type & Primary channel & Cross-world wedge (headline object) \\
\midrule
Uniform / hub / fringe transport & 1, 3, 4 & $\Delta\Pi^{- \theta}$, OD flows: \textbf{larger} \\
Hub productivity only & local $\bar A$ & Population / OD: \textbf{negligible} \\
Outer-ring amenity only & local $\bar u$ & Population / OD: \textbf{negligible} \\
Directional capacity (Stage~2) & 1 + congestion & Not identified in symmetric world \\
\bottomrule
\end{tabular}
\end{center}

This taxonomy motivates Stage~2 tidal-lane analysis: directionally targeted capacity
reallocation is a transport-technology shock that requires a directed calibration world;
symmetric measurement error is not innocuous even for \emph{non-directional} transport
policies evaluated on the wrong baseline.

\subsection{Mapping model objects to real planning outcomes}\label{sec:outcomes_map}

Stage~1 does not yet calibrate school seats, hospital beds, or transit line loads
directly.
The following mapping is the bridge from model outputs to policies planners discuss;
numbers in Table~\ref{tab:outcomes} come from the solved equilibrium in each world.

\begin{table}[htbp]
  \centering
  \caption{From model objects to planning outcomes (directed world, policy deltas)}
  \label{tab:outcomes}
  \small
  \begin{tabular}{p{0.22\textwidth}p{0.30\textwidth}rr}
  \toprule
  Model object & Planning interpretation & Hub access & Fringe transport \\
  \midrule
  $\Delta$ outer-ring $l_i^R$ (pp) & Residential public-service catchment & $-0.30$ & $+0.32$ \\
  $\Delta$ hub-bound OD share (pp) & Peak transit demand to CBD / hubs & $+0.78$ & $-0.25$ \\
  $\Delta$ hub $l_j^F$ (pp) & Employment-centre service / retail load & $+0.78$ & $-0.25$ \\
  OD dissimilarity (\%) & Commute-pattern / corridor investment priority & 6.6 & 0.55 \\
  p95 $|\Delta\log\Pi_i^{-\theta}|$ & Upper-tail job-access inequality & 0.025 & 0.032 \\
  \bottomrule
  \end{tabular}
  \vspace{0.4em}

  \begin{minipage}{0.92\textwidth}
  \footnotesize
  \emph{Notes:} Deltas are within the directed calibration world relative to that
  world's pre-shock equilibrium (+15\% hub-edge speed; $-20\%$ fringe-edge speed).
  ``Outer-ring'' is the top quartile of distance from the employment-weighted grid
  centroid ($\approx 25\%$ of nodes). Hub set: top employment quartile by observed
  $L_j^F$. Symmetric-world \emph{within-run} deltas are almost identical on population
  and OD shares; cross-world gaps after the same-named shock are an order of magnitude
  smaller on shares ($\lesssim 0.02$ pp on hub-bound commuting) but reach
  $+3.7\%$ on p95 $\Pi^{- \theta}$ for hub-access CF (Table~\ref{tab:policy_outcomes}).
  \end{minipage}
\end{table}

\paragraph{Residential public services.}
Grid-level residence shares $L_i^R=\Lbar\,l_i^R$ scale demand for schools, primary care,
and local retail.
Outer-ring amenity (+10\%) raises fringe $l_i^R$ by $+0.25$ pp; fringe transport
(+20\% speed on outer edges) raises it by $+0.32$ pp, while hub-access improvements
draw $-0.30$ pp out of the fringe---the same shock in the symmetric world moves shares
by virtually the same amount because the shock rule is spatial, not directional.

\paragraph{Employment-centre and corridor services.}
Predicted OD flows $L_{ij}$ identify origin--destination bundles that load on radial
arteries and metro lines.
Hub-access CF reweights $\approx 6.6\%$ of bilateral commuter mass (OD dissimilarity)
and raises the share of commuters with hub workplaces by $+0.78$ pp---a directly
observable object for peak transit scheduling and parking supply near employment centres.

\paragraph{Access inequality.}
Even when population shares barely differ across calibration worlds, upper-tail job
access $\Pi_i^{-\theta}$ does: baseline p95 $|\Delta\log\Pi_i^{-\theta}|\approx 0.52$
across worlds (Figure~\ref{fig:ma_maps}), widening to $+3.7\%$ \emph{relative} p95 gap
after hub-access CF.
Policies that target ``left-behind'' residential zones should be evaluated on
$\Pi_i^{-\theta}$ maps, not on $\Wbar$ alone.

\paragraph{What directional congestion does \emph{not} move in Stage~1.}
Pure scalings of hub $\bar A_j$ or fringe $\bar u_i$ produce $<0.01$ pp cross-world
wedges on both $\Wbar$ responses and population shares: the geographic shock sets are
identical and fundamentals are held fixed within each world.
Directional bias bites when the policy changes \emph{how} commuters reach places
through $T_{ij}$.

\subsection{Non-directional policy counterfactuals}\label{sec:policy_cf}

To separate \emph{directionality of measurement} from \emph{directionality of policy},
we apply five fixed-cost shocks that do \emph{not} reallocate lanes across directions:
\begin{enumerate}[leftmargin=2em,itemsep=0.2em]
  \item \textbf{Uniform speed (+10\%):} $\tau_{kl}\leftarrow 0.9\,\tau_{kl}$ everywhere.
  \item \textbf{Hub accessibility (+15\% speed on hub edges):} factor $0.85$ on edges
        touching the top employment quartile by observed $L_j^F$.
  \item \textbf{Hub productivity (+10\%):} $\bar A_j\leftarrow 1.1\,\bar A_j$ on the same
        hub set; travel costs unchanged.
  \item \textbf{Outer-ring amenity (+10\%):} $\bar u_i\leftarrow 1.1\,\bar u_i$ on fringe
        cells in the top quartile of distance from the employment-weighted grid centroid
        ($\approx 25\%$ of nodes).
  \item \textbf{Fringe transport (+20\% speed on fringe edges):} factor $0.80$ on edges
        touching the same outer-ring set.
\end{enumerate}
In each calibration world, $(\bar u,\bar A,\chi)$ are held at that world's inverted
values; only $\tau$ and/or hub $\bar A_j$ change; the static system is re-solved in levels
(no exact-hat, no congestion feedback).

Table~\ref{tab:policy_cf} and Figure~\ref{fig:policy_cf}.
Transport shocks move residence sorting and OD structure within each world; the
largest \emph{cross-world} gaps appear on $\Wbar$ responses (hub access: 0.14 pp
wedge) and on upper-tail job access after hub-access CF ($+3.7\%$ relative p95
$\Pi^{- \theta}$ gap; Table~\ref{tab:policy_outcomes}).
Pure local shocks to inverted fundamentals---hub productivity and outer-ring
amenity---produce negligible cross-world wedges on both welfare and observable
population/OD outcomes ($\lesssim 0.01$ pp on shares), confirming that calibration-world
bias propagates when policy loads on mis-measured commuting technology (Channels 1
and 4), not when it only scales local $\bar A_j$ or $\bar u_i$ on fixed geographic sets.

\begin{table}[htbp]
  \centering
  \caption{Cross-world outcome wedges after the same-named policy shock}
  \label{tab:policy_outcomes}
  \small
  \begin{tabular}{lrrrr}
  \toprule
  Policy & Hub-bound OD gap (pp) & Outer-ring $l^R$ gap (pp) & p95 $\Pi^{- \theta}$ gap (\%) & $\Wbar$ wedge (pp) \\
  \midrule
  Uniform speed (+10\%)      & $-0.01$ & $+0.01$ & $+0.20$ & $-0.07$ \\
  Hub access (+15\% on hubs) & $-0.02$ & $+0.01$ & $+3.72$ & $-0.14$ \\
  Hub productivity (+10\%)   & $-0.001$ & $+0.000$ & $-4.05$ & $-0.006$ \\
  Outer-ring amenity (+10\%) & $+0.000$ & $-0.000$ & $-4.19$ & $+0.006$ \\
  Fringe transport (+20\%)   & $-0.001$ & $+0.001$ & $-3.15$ & $+0.005$ \\
  \bottomrule
  \end{tabular}
  \vspace{0.4em}

  \begin{minipage}{0.92\textwidth}
  \footnotesize
  \emph{Notes:} Each column is symmetric-world minus directed-world \emph{after} applying
  the same policy rule in both worlds (post-shock outcomes, not response deltas).
  Hub-bound OD share $=\sum_{i,j\in\mathrm{hub}} L_{ij}/\sum_{i,j}L_{ij}$.
  p95 gap is percent difference in p95 $\Pi_i^{-\theta}$ across worlds.
  Population-share gaps stay tiny because re-inversion pins marginals at calibration;
  job-access tails and $\Wbar$ responses carry the cross-world transport bias.
  \end{minipage}
\end{table}

\begin{table}[htbp]
  \centering
  \caption{Non-directional policy CFs: directed vs.\ symmetric calibration world}
  \label{tab:policy_cf}
  \small
  \begin{tabular}{lrrrrrr}
  \toprule
  Policy & $\Delta\Wbar$ dir (\%) & $\Delta\Wbar$ sym (\%) & Wedge (pp)
  & Res.\ dissim.\ dir (\%) & OD dissim.\ dir (\%) & OD dissim.\ sym (\%) \\
  \midrule
  Uniform speed (+10\%)        & 0.823 & 0.755 & $-0.068$ & 1.81 & 4.60 & 4.68 \\
  Hub access (+15\% on hubs)   & 0.167 & 0.025 & $-0.142$ & 2.86 & 6.57 & 6.67 \\
  Hub productivity (+10\%)     & $-0.692$ & $-0.698$ & $-0.006$ & 0.38 & 0.45 & 0.45 \\
  Outer-ring amenity (+10\%)   & 0.351 & 0.357 & $+0.006$ & 0.25 & 0.25 & 0.25 \\
  Fringe transport (+20\%)   & 0.585 & 0.590 & $+0.005$ & 0.35 & 0.55 & 0.55 \\
  \bottomrule
  \end{tabular}
  \vspace{0.4em}

  \begin{minipage}{0.92\textwidth}
  \footnotesize
  \emph{Notes:} Wedge $=\Delta\Wbar^{sym}-\Delta\Wbar^{dir}$.
  Dissimilarities are $0.5\sum|\Delta\text{share}|$ relative to each world's pre-shock
  equilibrium. Hub set: top quartile of $L_j^F$. Outer-ring set: top quartile of
  distance from the employment-weighted grid centroid among residential cells.
  \end{minipage}
\end{table}

\begin{figure}[htbp]
  \centering
  \begin{minipage}{0.48\textwidth}
    \centering
    \includegraphics[width=\linewidth]{policy_cf_by_world.png}
    \caption{Policy responses by calibration world}\label{fig:policy_cf}
  \end{minipage}\hfill
  \begin{minipage}{0.48\textwidth}
    \centering
    \includegraphics[width=\linewidth]{policy_cf_welfare_wedge.png}
    \caption{Welfare-response wedge (sym $-$ dir)}\label{fig:policy_wedge}
  \end{minipage}
\end{figure}

\subsection{Asymmetry simulations}\label{sec:simulations}

Each simulation point re-inverts fundamentals at the stated cost surface.

\textbf{Intensity} $\rho\in\{0,0.25,\ldots,1\}$:
$\tau_{kl}(\rho)=\tau^{\sym}_{kl}\exp\{\rho[\log\tau^{\obs}_{kl}-\log\tau^{\sym}_{kl}]\}$.
Moving from $\rho=0$ (symmetric) to $\rho=1$ (directed) raises $\Wbar$ by $1.15\%$ while
leaving residence shares essentially unchanged ($<0.001\%$ dissimilarity at $\rho=1$).

\textbf{Threshold} $c\in\{1,1.05,\ldots,2\}$:
keep $\tau^{\obs}$ only on reciprocal pairs with $r_{kl}\ge c$, else symmetrise; compare
to the fully symmetric world.
At $c=2$ (8.9\% of reciprocal edges remain directed), $\Wbar$ still rises by $0.86\%$
relative to the symmetric calibration.

\begin{figure}[htbp]
  \centering
  \begin{minipage}{0.48\textwidth}
    \centering
    \includegraphics[width=\linewidth]{simulation_asymmetry_intensity.png}
    \caption{Asymmetry intensity $\rho$}\label{fig:intensity}
  \end{minipage}\hfill
  \begin{minipage}{0.48\textwidth}
    \centering
    \includegraphics[width=\linewidth]{simulation_asymmetry_threshold.png}
    \caption{Asymmetry threshold $c$}\label{fig:threshold}
  \end{minipage}
\end{figure}

\subsection{Tidal-lane prototype: short-run vs.\ long-run}\label{sec:tidal_lane}

We add a first-pass tidal-lane exercise on the directed calibration world while
keeping Stage~1 scope: \emph{no congestion propagation} to untreated links.
For each selected reciprocal pair, one lane is reallocated from the fast direction
to the slow direction (BPR at fixed observed congestion $\xi$, $\delta_1=0.488$), or
--- in a sensitivity arm --- peak times are scaled directly ($-10\%$ slow,
$+10\%$ fast).
Pairs are ranked by peak travel-time ratio $r_{kl}$; filters require length
$>200$\,m and at least two lanes on the fast direction.

\paragraph{Horizons.}
\begin{itemize}[leftmargin=2em,itemsep=0.2em]
  \item \textbf{Short-run:} fix $(l_i^R,l_j^F)$ at the baseline equilibrium; recompute
        $\Pi_i$, $\Wbar$, and implied OD flows through the new $T_{ij}$.
  \item \textbf{Long-run:} re-solve \eqref{eq:eq_R}--\eqref{eq:eq_F} with sorting.
\end{itemize}

Runner: \texttt{run\_tidal\_lane\_counterfactuals.py}; outputs under
\texttt{counterfactuals/tidal\_lane/}.
The default grid varies top-$K\in\{10,25,50,100\}$, $\Delta n\in\{0.5,1\}$ lanes,
minimum asymmetry $r\ge\{1.10,1.25,1.50\}$, plus three direct-$\tau$ scale bundles.

Table~\ref{tab:tidal_lane} and Figure~\ref{fig:tidal_short_long} summarise the
central pattern.
\textbf{Short-run} access gains are small but positive (up to $\approx +0.08\%$ on
$\Wbar$ for the largest direct-$\tau$ bundle on 50 pairs; BPR top-100:
$+0.075\%$), with implied OD reallocation up to $\approx 1\%$ dissimilarity.
\textbf{Long-run} sorting \emph{reverses} the sign: the same BPR shocks reduce
$\Wbar$ once population and workplaces re-equilibrate (e.g.\ top-50, $\Delta n=1$:
$-0.07\%$; residence sorting $\approx 0.35\%$; OD $\approx 0.6\%$).
Intuitively, treating a handful of adjacent grid links ignores network-wide rerouting
and general-equilibrium feedback; the exercise nonetheless shows that directional
capacity reallocation is not well approximated by a symmetric average cost, and
that short-run access calculations overstate net gains if residential sorting is
suppressed.

\begin{table}[htbp]
  \centering
  \caption{Tidal-lane prototype (directed world; no congestion spillovers)}
  \label{tab:tidal_lane}
  \small
  \begin{tabular}{llrrrrr}
  \toprule
  Spec & Horizon & Pairs & $\Delta\Wbar$ (\%) & Res.\ dissim.\ (\%) & OD dissim.\ (\%) & Mean $\Delta\tau$ slow (\%) \\
  \midrule
  BPR top-100, $\Delta n=1$ & Short & 100 & $+0.075$ & 0 & 0.99 & $-12.1$ \\
  BPR top-100, $\Delta n=1$ & Long  & 100 & $-0.140$ & 0.52 & 0.93 & $-12.1$ \\
  BPR top-50, $\Delta n=1$  & Short & 50  & $+0.037$ & 0 & 0.62 & $-12.2$ \\
  BPR top-50, $\Delta n=1$  & Long  & 50  & $-0.070$ & 0.35 & 0.60 & $-12.2$ \\
  Direct $\tau$: $-15\%/+15\%$, top-50 & Short & 50 & $+0.079$ & 0 & 0.71 & $-15.0$ \\
  Direct $\tau$: $-15\%/+15\%$, top-50 & Long  & 50 & $-0.085$ & 0.43 & 0.73 & $-15.0$ \\
  \bottomrule
  \end{tabular}
  \vspace{0.4em}

  \begin{minipage}{0.92\textwidth}
  \footnotesize
  \emph{Notes:} BPR uses fixed $\xi$ at observed flows; fast direction slows because
  lanes fall. Long-run rows re-solve sorting; short-run rows fix baseline shares.
  Effects are an order of magnitude smaller than hub-access CFs
  (\S\ref{sec:policy_cf}) because only 50--100 adjacent pairs are treated on a
  2642-edge network.
  \end{minipage}
\end{table}

\begin{figure}[htbp]
  \centering
  \includegraphics[width=0.82\textwidth]{tidal_lane_short_vs_long.png}
  \caption{Tidal-lane grid: short-run vs.\ long-run $\Wbar$ and OD responses}
  \label{fig:tidal_short_long}
\end{figure}

\subsection{Interpretation and story}\label{sec:interpretation}

Stage~1 delivers a coherent measurement narrative rather than a welfare ranking:
\begin{enumerate}[leftmargin=2em,itemsep=0.35em]
  \item \textbf{Directional congestion is widespread} in the retained Beijing network
        (median peak/off-peak ratio 1.23 on reciprocal links; Figure~\ref{fig:asym_hist}).

  \item \textbf{Symmetric link indices mis-state commuting technology} even when both
        calibration worlds fit the same residence and workplace marginals.
        Predicted OD flows differ by $\approx 2.5\%$; grid-level job access
        $\Pi_i^{-\theta}$ differs at p95 $|\Delta\log|\approx 0.52$
        (Figure~\ref{fig:ma_maps}); inverted fundamentals compensate
        (p95 $|\log(\bar u_i^{\sym}/\bar u_i^{\dir})|\approx 0.55$).

  \item \textbf{$\Wbar$ is not the right headline} under re-inversion: baseline
        $\Wbar$ falls 1.14\% in the symmetric world mainly because $(\bar u,\bar A)$
        re-adjust, not because observed population counts change.
        We report $\Wbar$ for completeness and because transport CFs still show
        asymmetric \emph{responses} (hub-access wedge 0.14 pp), but planning-relevant
        objects are population shares, OD corridors, and $\Pi_i^{-\theta}$ maps
        (\S\ref{sec:outcomes_map}).

  \item \textbf{Which policies care about directional measurement?}
        \emph{Transport} shocks---uniform speed-up, hub-edge accessibility, fringe
        connectors---reallocate OD mass ($4.6$--$6.7\%$ dissimilarity) and fringe/hub
        population within each world, and widen cross-world upper-tail job-access gaps
        (hub access: $+3.7\%$ on p95 $\Pi^{- \theta}$).
        \emph{Place-based} shocks that only scale hub $\bar A_j$ or fringe $\bar u_i$
        on fixed zones produce identical population and OD responses across worlds
        ($<0.02$ pp share gaps); they are poor diagnostics of directional bias in
        Stage~1.

  \item \textbf{Tidal-lane prototype (\S\ref{sec:tidal_lane}).}
        Directional lane reallocation without congestion spillovers produces modest
        short-run access gains but negative long-run $\Wbar$ once sorting adjusts;
        implied OD shifts ($\sim$0.6--1\%) remain far below hub-access transport CFs
        because only top asymmetric \emph{adjacent} pairs are treated.
        Stage~2 should combine directional shocks with network-wide propagation.
\end{enumerate}

Stage~1 does \emph{not} support full tidal-lane welfare claims or structural $\lambda$
estimation; the prototype in \S\ref{sec:tidal_lane} is deliberately partial (no
congestion spillovers, adjacent-pair sample only).
It \emph{does} motivate Stage~2: policies that explicitly reallocate directed capacity
(tidal lanes) cannot be evaluated on a symmetric-cost calibration that already hides
peak-direction congestion and shifts $(\bar u,\bar A)$ to compensate.
The empirical next step is to export zone-level $\Delta L_i^R$, $\Delta L_{ij}$, and
$\Delta\Pi_i^{-\theta}$ from directed CFs into service-planning metrics (school
catchments, transit peak load, hospital referral zones)---Stage~1 supplies the spatial
equilibrium inputs; the institutional mapping is deliberately left for the applied
section of the paper.

\section{Link to Stage~2}\label{sec:stage2}

Stage~2 activates congestion, anchors travel times to observed speeds, allows
congestion spillovers across the network, and evaluates lane-capacity counterfactuals
(including tidal lanes on corridor samples) under fundamentals recovered here.
Stage~1 supplies: (i)~directed vs symmetric measurement bias quantification;
(ii)~policy taxonomy for which outcomes load on $T_{ij}$; (iii)~tidal-lane prototype
code path (\texttt{tidal\_lane\_counterfactuals.py}) to be extended with endogenous
$\xi$ updates.

\appendix
\section{Module map (code)}\label{app:code}

\begin{center}
\small
\begin{tabular}{p{0.42\textwidth}p{0.50\textwidth}}
\toprule
File & Role \\
\midrule
\texttt{full\_estimation/build\_primitives.py} & Load QSM parquet $\to$ \texttt{Primitives} \\
\texttt{full\_estimation/invert\_fundamentals.py} & Fixed-point inversion \eqref{eq:invert_R}--\eqref{eq:invert_F} \\
\texttt{full\_estimation/parameters.py} & Frozen $(\theta,\alpha,\beta)$ \\
\texttt{directional\_cost/fixed\_tau\_model.py} & $T$, equilibrium, OD gravity, comparisons \\
\texttt{directional\_cost/run\_directional\_cost.py} & Main Stage~1 driver \\
\texttt{directional\_cost/policy\_counterfactuals.py} & Non-directional CF specs \\
\texttt{directional\_cost/tidal\_lane\_counterfactuals.py} & Pair selection, BPR shock, horizons \\
\texttt{directional\_cost/od\_fit\_estimation.py} & Optional $\theta$/SMM diagnostics (parallel track) \\
\bottomrule
\end{tabular}
\end{center}

Detailed session log:
\path{Documents/L2_notes/stage1_fixed_directed_cost_complete_handoff_20260628.md}.

\bibliographystyle{aer}
\bibliography{tidal_lanes_refs}

\end{document}
